\part{初高衔接}
\chapter{式与方程}

\section{因式分解}
\subsection{十字相乘，立方和与差及其拓展}
\begin{know}
\item \textbf{十字相乘}

能用十字相乘法分解一元二次式,
能通过选取合适的主元来分解具有多个变量的代数式.

\item \textbf{立方和与差公式}

熟记立方和与差公式以及更高次幂的拓展, 能结合拆项补项的技巧分解一元三次式, 
乃至于特殊情况的一元高次式.
\end{know}

所谓因式分解\index{因式分解}，即把一个多项式写成几个多项式乘积的形式。在初中我们学习过用提公因式法，公式法来因式分解。
而在这里我们将介绍更多的公式与常见的方法。

首先是十字相乘\index{因式分解!十字相乘}，这针对的是形如$x^2+px+q$或$ax^2+bx+c$的形式二次三项式。由于有乘法公式
\[(x+m)(x+n)=x^2+(m+n)x+mn\]
\[(px+m)(qx+n)=pqx^2+(pn+qm)x+mn\]

我们不难想到，如果要把式子从右到左写出来，就要分解二次三项式的常数项（和二次项），并使他们交叉相乘后为一次项系数。
显然这个方法不乏一些运气成分，但是并非毫无规律可循的。先分解质因数，再调整正负号，得出来结果并不难。
\begin{example}
用十字相乘法分解因式
 $(1)x^2+4x-5$

 $(2)2x^2+x-1$
\end{example}
\begin{jie}
    
\[(1)x^2+4x-5=(x-5)(x+1)\]
\[(2)2x^2+x-1=(2x-1)(x+1)\]
\qed 
\end{jie}

\begin{exercise}
\label{exer:exercise1.1.1}
    分解因式
\par$(1)x^2-4x-12$\qquad\ \ \ \ \ \ \ \ \,$(2)x^2-8x-15$
\par$(3)x^2-3x-10$\qquad\ \ \ \ \ \ \ \ \,$(4)x^3-4x^2-21x$
\par$(5)x^2+12x+32$\qquad\ \ \ \ \ \ \ $(6)x^4+6x^2+8$
\par$(7)a^2-4a-21$\qquad\ \ \ \ \ \ \ \, $(8)y^2-10x+16$
\par$(9)6x^2+13x+6$\qquad\ \ \ \ \ \ \ $(10)7x^2-19x-6$
\par$(11)12x^2-13x+3$\qquad\ \ \ \ $(12)2x^2-5x-12$
\par$(13)4k^2+4k-15$\qquad\ \ \ \ \ \,$(14)15m^2+m-2$

\end{exercise}

在过去，我们学习了完全平方公式\index{因式分解!完全平方公式}和平方差公式\index{因式分解!平方差公式}，即
\[(a\pm b)^2=a^2\pm 2ab +b^2\]
\[(a+b)(a-b)=a^2-b^2\]

我们如何将它们变得更难呢？

一种是变元数量的增加，比如将$(a+b)^2$升级为$(a+b+c)^2$,甚至$(x_1+x_2+\cdots +x_n)^2$

显然有\[(a+b+c)^2=a^2+b^2+c^2+2ab+2bc+2ac\]
\[(\sum_{k= 1}^{n}x_k  )^2=\sum_{\substack{1\leq i\leq n\\1\leq j\leq n}}x_ix_j\]

另一种是次数的增加，比如将$(a+b)^2$升级为$(a+b)^n$，将$a^2-b^2$升级为$a^n-b^n$。
前者的结果需要排列组合的背景，即以后学到的二项式定理。
而后者则是比较容易的，请验证下列式子并尝试自己归纳规律。
\begin{align*}
    a^2-b^2&=(a-b)(a+b)\\
    a^3-b^3&=(a-b)(a^2+ab+b^2)\\
    a^4-b^4&=(a-b)(a^3+a^2b+ab^2+b^3)\\
    a^5-b^5&=(a-b)(a^4+a^3b+a^2b^2+ab^3+b^4)\\
    \vdots  \\
    a^n-b^n&=(a-b)(a^{n-1}+a^{n-2}b+\cdots +ab^{n-2}+b^{n-1})\\
           &=(a-b)(\sum_{k=1}^{n}a^{n-k}b^{k-1})
\end{align*}

如果令$b=1$，还可以得到另一个常用的式子
\begin{align*}
     a^n-1&=(a-1)(a^{n-1}+a^{n-2}+\cdots +a+1)\\
           &=(a-1)(\sum_{k=1}^{n}a^{k-1})
\end{align*}
     
显然，用$\frac{a}{b}$代替$a$，就可以由$a^n-1$得到$a^n-b^n$。这些都是常用的变形技巧。

对于$n$为奇数的式子，用$-b$代替$b$，就可以轻松得到立方和乃至更高奇数次方差的公式。
\begin{align*}
a^3+b^3&=(a+b)(a^2-ab+b^2)\\
a^5+b^5&=(a+b)(a^4-a^3b+a^2b^2-ab^3+b^4)\\
\vdots \\
a^n+b^n&=(a+b)(a^{n-1}-a^{n-2}b+\cdots -ab^{n-2}+b^{n-1})\\
       &=(a+b)(\sum_{k=1}^{n}a^{n-k}(-b)^{k-1})\qquad(n=2k+1,k\in \mathbb{N} )
\end{align*}

你要怎么记忆这两组式子呢？请尝试着找出方便记忆的规律。

当然，特别高次的式子并不常用，我们还是更关心不那么难的立方和与差公式，请先完成下面几个简单的因式分解为下面更复杂的做准备。
\begin{exercise}
    用立方和与差公式分解因式
    \begin{align*}
        (1)&x^3-1\\
        (2)&64x^3+8\\
        (3)&x^3-27
    \end{align*}
\end{exercise}

这几个例子还是很简单的对吧，但是如果我们面对的是形如$ax^3+bx^2+cx+d$的三次四项式
（为了使难度不那么高，往往会少一个二次项或一次项）呢？

面对这样的一元三次式的因式分解，我们常常使用的是拆项补项的技巧。值得说明的是这种技巧的应用相当广泛，我们会不断地见到它。

\begin{example}
用拆项补项的方法分解因式。  
\begin{align*}
    x^3-7x+6&=x^3-x-(6x-6)\\
            &=x(x^2-1)-6(x-1)\\
            &=x(x+1)(x-1)-6(x-1)\\ 
            &=(x-1)(x^2+x-6)\\  
            &=(x-1)(x+3)(x-2)
\end{align*}%对齐公式

另一种选择是：
\begin{align*}
    x^3-7x+6&=x^3-7x+(7-1)\\
            &=x^3-1-7x+7\\
            &=(x-1)(x^2+x+1)-7(x-1)\\
            &=(x-1)(x^2+x-6)\\  
            &=(x-1)(x+3)(x-2)
\end{align*}
\end{example}

\begin{exercise}
    用拆项补项的方法分解因式。  

    %\begin{enumerate}[(1)]
    %    \item a
    %    \item b
    %\end{enumerate}

$(1)x^3+x^2-36$

$(2)x^3+3x-4$

$(3)y^3-6y-4$

$(4)4x^3-31x+15$

$(5)x^3-2x^2-4x+8$

$(6)m^4+4$
\end{exercise}

在这节的最后，提出两个问题：

究竟为什么要因式分解？答案很明显，我们要更快更方便地解方程（至少在高中阶段是这样的）。

什么时候可以用因式分解，我拿到这个式子怎么就知道可以用上述方法去解开呢？这个可以稍微思考一下。

这节的标题中出现了一个陌生的词汇，主元。元，说白了就是变量，主元就是起主导地位的自变量。
因为一个题里面大概率不会只有一个字母$x$，还可能有其他$a,b,m,n,$……虽然把它们当成同等地位的变量处理多数也可以做。
但是若使用主元法的想法，把一个字母看做主元（主要的自变量），其他字母看成参数，在处理起来会有事半功倍的效果。

我们在这一节仅仅是大概看一下再因式分解中的用途，便于找变量之间的关系，具体方法看下例。
\begin{example}
\[3b^2+4k^2+8kb-2b-1\]
\end{example}
\begin{jie}
不妨以$k$为主元，$b$为参数。此时是关于\(k\)的二次三项式。 
\begin{align*}
    3b^2+4k^2+8kb-2b-1&=4k^2+8bk+(3b^2-2b-1)
\\                    &=4k^2+8bk+(b-1)(3b+1)
\\                    &=(2k+b-1)(2k+3b+1)
\end{align*}
\end{jie}
\begin{exercise}\label{exer:exercise 1.2.1}
    选取合适的主元，分解下列因式
    
$(1)x^2-4mx+8mn-4n^2$

$(2)6x^2-7xy+2y^2$

$(3)3x^2+4xy-4y^2+8x-8y-3$
\end{exercise}
\begin{zhu}
你能选取别的主元再做一遍上面的题目吗？如果不用主元法的方法思考可以做吗？
\end{zhu}
\begin{example}选取合适的主元，分解下列因式：

    \((1)a^2c-a^2b-bc^2-b^3\);

    \((2)(x_0^2+y_0^2-1)^2=12y_0^2+16(x_0+1)^2\).
\end{example}
\begin{analyse}
    在上几个例子中，待分解的次数都不高，因此选择变量作为主元并不是一个特别需要思量的事情。但是这两个次数都偏高，尤其$(2)$
    展开后应该最高出现了$4$次项，故应该仔细观察选择哪个字母作为主元可以使问题尽可能化简。

    \((1)\)中，\(b\)出现了\(3\)次项，首先被排除，毕竟二次项的分解我们更加熟悉；选择\(a\)
    或\(c\)作为主元应该差不多，毕竟都是二次项。

    \((2)\)中，理论上只有\(x_0,y_0\)和常数。如果选择\(x_0\)作为主元，其次数有\(4,2,1\)，这显然不太好处理；但是
    选择\(y_0\)的话就可以发现其次数只有\(4,2\)，换句话说可以将\(y_0^2\)这个整体看做主元，得到一个二次三项式。
\end{analyse}
\begin{jie}
    \((1)\)

    \((2)\)先将式子写为$[y_0^2+(x_0^2-1)]^2-12y_0^2-16(x_0+1)$, 将\(y_0^2\)看做主元。

    依次计算二次项系数\(: 1\), 
    
    一次项系数\(: 2(x^2_0-1)-1=2(x_0^2-7)\), 
    
    常数项\(:  (x_0^2-1)^2-16(x_0+1)^2=(x_0+1)^2(x_0-1)^2-16(x_0+1)^2=(x_0+1)^2[(x_0-1)^2-16]=(x_0+1)^2(x^2-2x_0-15)\)

    故原式可以写作$y_{0}^{4}+\left(2 x_{0}^{2}-14\right) y_{0}^{2}+\left(x_{0}^{2}+2 x_{0}+1\right)\left(x_{0}^{2}-2 x_{0}-15\right)=0$

    试图使用十字相乘的想法，会发现恰好有  $\left(x_{0}^{2}+2 x_{0}+1\right)+\left(x_{0}^{2}-2 x_{0}-15\right)=2 x_{0}^{2}-14  $，因此可以将条件式因式分解为
 $$   \left(y_{0}^{2}+x_{0}^{2}+2 x_{0}+1\right)\left(y_{0}^{2}+x_{0}^{2}-2 x_{0}-15\right)=0$$ .
\qed
\end{jie}

\clearpage
\subsection{更高次幂的因式分解与长除法*}
\begin{know}
    \item \textbf{一元多项式根的性质}
    
    \item \textbf{长除法}
\end{know}
在这一节要简单地回答一下1.1节中提出的问题，什么时候可以因式分解？
真的需要分解一些高次式的时候，有没有什么比较初等的思路？

对于一个二次式而言，答案是显然的。我们只要计算二次式$ax^2+bx+c$的判别式$\Delta =b^2-4ac$即可。
如果是完全平方数或者完全平方式，那么就可以因式分解。
\begin{comment}
显然这里的因式分解指的是在有理数域上因式分解，我们并不希望式子的系数出现根号。
\end{comment}
我们用上两节的题目举个例子。
\begin{example}
    计算下列式子的$\Delta$

    $(1)12x^2-13x+3$\qquad $(\text{练习}\ref{exer:exercise1.1.1}(11))$

    $(2)3x^2+4xy-4y^2+8x-8y-3$\qquad$(\text{练习} \ref*{exer:exercise 1.2.1}(3))$
\end{example}

\begin{jie}
    $(1)$\[\Delta=(-13)^2-4\times12\times3=25\]

    $(2)$以$x$为主元整理式子为
\[3x^2+(4y+8)x-(4y^2+8y+3)\]
\begin{align*}
    \Delta&=(4y+8)^2+4\times3(4y^2+8y+3)\\
          &=4(2y+4)^2+4(12y^2+24y+9)\\
          &=4(4y^2+16y+16+12y^2+24y+9)\\
          &=4(16y^2+40y+25)\\
          &=4(4y+5)^2
\end{align*}
\end{jie}

其实我们就这样接着套求根公式的话，也能很快得出结果。这两种方法我个人认为很难评优劣
如果真的在考场上遇到了，似乎花费的时间也不会差很多，更多还是看个人的计算习惯。

但是我们遇到了三次式呢？我们并不会三次方程的求根公式和判别式，
因此我们需要一些小技巧来猜出来分解后的结构。
在此，不加证明地给出几条定理。
\begin{comment}
    下面的定理为了追求严谨，并非高中生能看懂的人话。请试着在旁边写下自己的翻译。
    在了解复数并定义多项式的一些基本知识后，证明是容易的。
\end{comment}
\begin{theorem}
   $ c\in \mathbb{Q}$ 为$f(x) \in \mathbb{Q} (x)$的根$\Leftrightarrow (x-c)|f(x)$
\end{theorem}
\begin{theorem}
    实数域上的不可约多项式是一次多项式，或者是含一对共轭非实复根的二次多项式。因而每个次数
    大于零的实系数多项式都可以分解为实数域上的一次或二次不可约多项式的乘积。
\end{theorem}
\begin{theorem}
若$f(x)\in \mathbb{Z} (x)$，且$f(x)=a_nx^n+a_{n-1}x^{n-1}+\cdots+a_1x+a_0(a_n,a_0\neq 0)$，
设有理根为既约分数$\frac{p}{q}$的形式，则必然有$p|a_0,q|a_n$成立。特别地，当$a_n=1$的时候，其有理根必然为其常数项$a_0$的因数。
\end{theorem}

尽管你一时半会理解不了它们，也不用担心。它们只是出于严谨放在那里的。
至少它们（高中所用的大多数定理都有如此优点）是符合数学直觉的且工整漂亮的。
下面我们开始正式操作一下面对高次式我们可以做的猜根和长除法。\index{长除法}
\begin{example}
    使用上述定理猜根，并用长除法分解
    \[2x^4+3x^3-16x^2+3x+2\]
\end{example}

\begin{jie}
\ 
\begin{way1}
    根据定理$1.3.3$，原多项式可能的根为$\pm1,\pm2,\pm\frac{1}{2}$。发现$x=2$是多项式的根。
    故$(x-2)$为原多项式的一个因式。
    
    由长除法得\[2x^4+3x^3-16x^2+3x+2=(x-2)(2x^3+7x^2-2x-1)\]
    
    而后面那一堆没有有理根了，故因式分解到这就算结束了。
\end{way1}
\begin{way2}
    根据定理$1.3.2$，原多项式必然可以写作两个二次式的乘积（无论它们是否可约）。
    不妨设\[2x^4+3x^3-16x^2+3x+2=(2x^2+ax+b)(x^2+cx+d)\]
    
    请注意这次原多项式系数$2$相当优秀，是一个质数，因此分解后的二次多项式的二次项系数直接确定下来了。
    接下来就是两侧系数对应相等的暴力计算了。在此打个预防针，根据法$1$可知这个数字大概率不会友好，也不一定唯一（为什么？）。

    右侧展开为
    \[2x^4+3x^3-16x^2+3x+2=2x^4+(a+2c)x^3+(ac+b+2d)x^2+(ad+bc)x+(bd)\]
    
    得到方程组
    \[  
    \left\{
    \begin{lgathered}
    a+2c=3\\
    ac+b+d=-16\\
    ad+bc=3\\
    bd=2
    \end{lgathered}
    \right.
    \]
    这种方程组的解法往往不统一，但主要思路还是消元降次，如果运气不好的话很可能最后结果归于解一个三次方程。
    
    经\sout{计算器}计算得到两组实数解\sout{（甚至还有三组复数解）}：
    \begin{equation*}
    \left\{
    {
    \begin{array}{*{20}{c}}
        {a=\frac{3-\sqrt{161}}{2}}\\ 
        {b=2\ \ \ \ \ \ \ }\\  
        {c=\frac{3+\sqrt{161}}{4}}\\
        {d=1\ \ \ \ \ \ \ }
    \end{array}
    }
    \right.,
    \left\{
    {
    \begin{array}{*{20}{c}}
        {a=\frac{3+\sqrt{161}}{2}}\\  
        {b=2\ \ \ \ \ \ \ }\\
        {c=\frac{3-\sqrt{161}}{4}}\\
        {d=1\ \ \ \ \ \ \ }
    \end{array}
    }
    \right.    
    \end{equation*}

    故\[2x^4+3x^3-16x^2+3x+2=(2x^2+\frac{3-\sqrt{161}}{2}x+2)(x^2+\frac{3+\sqrt{161}}{4}x+1)\]

    或\[2x^4+3x^3-16x^2+3x+2=(2x^2+\frac{3+\sqrt{161}}{2}x+2)(x^2+\frac{3-\sqrt{161}}{4}x+1)\]
\end{way2}
    请自行验证法$1$和法$2$的结果相同。但是说实话，这种待定系数法过于笨重，只要了解大概即可。
    在后面我们会经常用这种方法来解决系数的问题的。

    \qed
\end{jie}

\
\begin{comment}
使用\mbox{\LaTeX}打长除法实在过于痛苦，请参考板书。
\end{comment}
\begin{att}
使用长除法的时，将多项式按高次到低次排列，缺少的项要用零补齐。
\end{att}
\begin{exercise}
    用猜根和长除法分解下列因式

    $(1)2x^3+7x^2-2x-1$

    $(2)x^3+5x^2-x-5$
\end{exercise}
\begin{exercise}
    已知$x^2+x+6$是多项式$x^4-2x^3+ax^2+bx+a+b-1$的因式，求$ab$的值。
\end{exercise}

\clearpage

\hw

\begin{homework}因式分解
    \begin{enumerate}
        \item 1
        \item 2
        \item 1
        \item 2
        \item 2
        \item 2
        \item 2
        \item 2
    \end{enumerate}

\end{homework}
\begin{homework}
    因式分解

    \begin{minipage}{0.45\textwidth}
         \begin{enumerate}
        \item \(ab-5bc-2a^2+10ac\)
        \item $3x²+8xy-3y²+4x+22y-7$
        \item \(ab-ac+bc-b^2\)
        \item \(4x^2-6x-y^2-3y\)
        \item \(a^2b-a^2c+b^2c-b^3\)
        \item \(a^2-b^2+4b-4\)
        \item \(\)
        \item 1
        \item 1
    \end{enumerate}
    \end{minipage}
    \hfill
    \begin{minipage}{0.45\textwidth}
        \begin{enumerate}
            \addtocounter{enumi}{+9}   
            \item \(ab-5bc-2a^2+10ac\)
            \item $3x²+8xy-3y²+4x+22y-7$
            \item \(ab-ac+bc-b^2\)
            \item \(4x^2-6x-y^2-3y\)
            \item \(a^2b-a^2c+b^2c-b^3\)
            \item \(a^2-b^2+4b-4\)
            \item \(\)
            \item 1
            \item 1
        \end{enumerate}
    \end{minipage}
   

\end{homework}
\begin{homework}

        
\end{homework}
\clearpage